% Source: https://artofproblemsolving.com/community/c3h1866692_permutations
% Stable identity: aops-topic-1866692-q01
(a) When we come to the $k$th office, the first $k-1$ offices have already been given to $k-1$ different members. Nobody may hold two offices, so the $k$th office can go to any of the remaining $n-(k-1)=n-k+1$ members, whatever the earlier choices were. For $k=r$ this gives $n-r+1$ choices.

(b) By the multiplication principle the number of ways is
\[
n(n-1)(n-2)\cdots(n-r+1)=\dfrac{n(n-1)\cdots(n-r+1)\cdot(n-r)!}{(n-r)!}=\dfrac{n!}{(n-r)!}.
\]
There are $r$ factors, one per office, and the last is $n-r+1$. For $n=10$ and $r=3$ this is $10\cdot9\cdot8=\boxed{720}$.
